(1) 取特征曲线满足 $dx/dt=u$，在其上 $du/dt=u_t+(dx/dt)u_x=u_t+uu_x=0$，所以速度保持初值 $u_0(\xi)$。由 $x(0)=\xi$ 积分得到 $x=\xi+t u_0(\xi)$，并有隐式解 $u(x,t)=u_0(\xi)$，其中 $\xi$ 由前一关系反求。
(2) 对特征关系求 $\partial x/\partial\xi=1+t u_0'(\xi)$。通过链式法则 $u_x=u_0'(\xi)/[1+t u_0'(\xi)]$。当分母接近零时梯度绝对值变得极大，虽然 $u$ 沿每条特征仍有限，因此破裂表现为波形陡化而非速度幅度发散。
(3) 若 $m_0=\min_\xi u_0'(\xi)<0$，第一个特征交叉发生在 $1+t m_0=0$，所以破裂时间 $t_*=-1/m_0>0$。若所有初始梯度均非负，向前特征不会因此交叉，经典解可在这个机理下继续平滑传播。
(4) 将方程写成守恒律 $u_t+f(u)_x=0$，$f(u)=u^2/2$。跨移动不连续面应用质量守恒的Rankine–Hugoniot条件 $s(u_R-u_L)=f(u_R)-f(u_L)$，得激波速度 $s=(u_L+u_R)/2$。凸通量的熵允许激波应有 $u_L>u_R$，此时左侧特征比激波快、右侧特征比激波慢。
